% BEGIN VOL08-EXPANSION VIII08
\section{Supplementary worked examples}
\label{sec:viii08-pedagogy}

\begin{example}[Two degree-one attachments]\label{ex:viii08-ped-01}
Any two degree-one maps $S^{n-1}\to S^{n-1}$ are homotopic, so attaching an $n$-cell along either gives homotopy-equivalent adjunction spaces under the standard CW hypotheses.
\end{example}

\begin{example}[Changing a loop representative]\label{ex:viii08-ped-02}
If two loops in a graph are homotopic through based loops, attaching a 2-cell along either loop imposes the same homotopy-level relation.  The geometric routes can differ while the resulting CW homotopy type agrees.
\end{example}

\begin{example}[Constant attachments]\label{ex:viii08-ped-03}
Any two based constant attaching maps $S^{n-1}\to X$ with the same base point coincide; attaching along the constant map gives a wedge $X\vee S^n$ in the standard based CW setting.
\end{example}

\section{Graded supplementary exercises}
The following exercises progress from direct computation and construction to proof, hypothesis testing, application, and synthesis.

\subsection*{Standard computations and constructions}

\begin{exercise}[Attachment data]\label{exr:viii08-ped-01}
Which part of an $n$-cell attachment may be replaced by a homotopic map?
\end{exercise}
\begin{hint}
The map on $S^{n-1}$.
\end{hint}
\begin{solution}
The attaching map $f:S^{n-1}\to X$ can be replaced by a homotopic attaching map $g$.
\end{solution}

\begin{exercise}[Constant attachment]\label{exr:viii08-ped-02}
What is the homotopy type of attaching an $n$-cell by a constant map to a based space $X$?
\end{exercise}
\begin{hint}
The boundary collapses to the base point.
\end{hint}
\begin{solution}
It gives $X\vee S^n$ (indeed the quotient has this standard form).
\end{solution}

\begin{exercise}[Degree classes]\label{exr:viii08-ped-03}
When $X=S^{n-1}$, what invariant classifies attaching maps $S^{n-1}\to S^{n-1}$ up to homotopy?
\end{exercise}
\begin{hint}
Use degree.
\end{hint}
\begin{solution}
Their degree in $\mathbb Z$ classifies the homotopy class for sphere self-maps.
\end{solution}

\begin{exercise}[Same degree]\label{exr:viii08-ped-04}
If two sphere self-maps have the same degree, what can be concluded about the corresponding one-cell attachments?
\end{exercise}
\begin{hint}
First identify their homotopy classes.
\end{hint}
\begin{solution}
They are homotopic, hence their adjunction spaces are homotopy equivalent under the standard attachment theorem.
\end{solution}

\begin{exercise}[Zero degree]\label{exr:viii08-ped-05}
What is special about a degree-zero attaching map $S^{n-1}\to S^{n-1}$?
\end{exercise}
\begin{hint}
It is null-homotopic in the sphere range.
\end{hint}
\begin{solution}
It is homotopic to a constant map, so the attachment has the homotopy type of $S^{n-1}\vee S^n$.
\end{solution}

\subsection*{Proofs}

\begin{exercise}[Homotopic attachment theorem]\label{exr:viii08-ped-06}
Outline why homotopic attaching maps give homotopy-equivalent adjunction spaces.
\end{exercise}
\begin{hint}
Use the attaching homotopy on a collar/cylinder.
\end{hint}
\begin{solution}
Insert the homotopy between the two boundary identifications into a mapping-cylinder comparison; the inclusion of the boundary sphere is a cofibration, allowing the homotopy to extend to homotopy equivalences of pushouts.
\end{solution}

\begin{exercise}[Degree consequence]\label{exr:viii08-ped-07}
Prove equal-degree sphere attachments yield homotopy-equivalent spaces.
\end{exercise}
\begin{hint}
Equal degree gives homotopic maps.
\end{hint}
\begin{solution}
Sphere self-maps with equal degree are homotopic; apply the homotopic attaching-map theorem.
\end{solution}

\begin{exercise}[Null attachment wedge]\label{exr:viii08-ped-08}
Prove a null-homotopic attaching map yields the homotopy type $X\vee S^n$.
\end{exercise}
\begin{hint}
Replace it by a constant attaching map.
\end{hint}
\begin{solution}
Homotopic attachments have the same homotopy type, and the constant attachment collapses $S^{n-1}$ to the base point, producing the wedge with $S^n$.
\end{solution}

\begin{exercise}[Iterated attachments]\label{exr:viii08-ped-09}
Explain how the theorem supports inductive comparisons of CW complexes.
\end{exercise}
\begin{hint}
Compare one skeleton at a time.
\end{hint}
\begin{solution}
If corresponding attaching maps are homotopic after identifying the lower skeleta, each attachment step preserves a homotopy equivalence, enabling cellular induction.
\end{solution}

\subsection*{Counterexamples and hypothesis tests}

\begin{exercise}[Homotopic gives homeomorphic]\label{exr:viii08-ped-10}
Test: homotopic attaching maps always give homeomorphic adjunction spaces.
\end{exercise}
\begin{hint}
The theorem guarantees homotopy type, not necessarily homeomorphism.
\end{hint}
\begin{solution}
False in that strength; homotopy equivalence is the robust conclusion.
\end{solution}

\begin{exercise}[Same image suffices]\label{exr:viii08-ped-11}
Test: attaching maps with the same image subset always yield the same homotopy type.
\end{exercise}
\begin{hint}
Winding/degree can differ despite identical image.
\end{hint}
\begin{solution}
False.  The homotopy class, not merely image, controls the attachment.
\end{solution}

\begin{exercise}[Different degree harmless]\label{exr:viii08-ped-12}
Test: changing the degree of a sphere attaching map never changes homology.
\end{exercise}
\begin{hint}
Degree becomes a cellular boundary coefficient.
\end{hint}
\begin{solution}
False.  Mapping-cone examples show different degrees produce different torsion homology.
\end{solution}

\subsection*{Applications and investigations}

\begin{exercise}[Model simplification]\label{exr:viii08-ped-13}
How can one simplify an attaching map before computation without changing homotopy type?
\end{exercise}
\begin{hint}
Replace it by a convenient representative of its homotopy class.
\end{hint}
\begin{solution}
Choose a homotopic map with simpler geometry or combinatorics; the resulting attached space retains the same homotopy type while calculations become easier.
\end{solution}

\begin{exercise}[Presentation robustness]\label{exr:viii08-ped-14}
Why does this theorem make CW descriptions robust under geometric perturbation?
\end{exercise}
\begin{hint}
Small/controlled deformations often preserve homotopy class.
\end{hint}
\begin{solution}
The exact geometric route of an attaching sphere can change while its homotopy class stays fixed, so homotopy-level invariants of the CW complex remain stable.
\end{solution}

\subsection*{Challenge problems}

\begin{exercise}[Degree-$m$ family]\label{exr:viii08-ped-15}
Compare the spaces obtained by attaching a 2-cell to $S^1$ with degrees $m$ and $-m$.
\end{exercise}
\begin{hint}
Their cellular homology depends on $|m|$; consider an orientation reversal of the cell or circle.
\end{hint}
\begin{solution}
They have isomorphic cellular homology and are in fact related by pre/postcomposition with degree $-1$ homeomorphisms, so the sign can be absorbed by orientation reversal; the magnitude records the essential torsion.
\end{solution}

\begin{exercise}[Attachment synthesis]\label{exr:viii08-ped-16}
Explain the hierarchy: equal maps, homotopic maps, equal degree maps on spheres, and resulting attached spaces.
\end{exercise}
\begin{hint}
Each implication weakens geometric data but preserves enough homotopy information.
\end{hint}
\begin{solution}
Equal maps are homotopic; on spheres equal degree implies homotopic; homotopic attaching maps yield homotopy-equivalent adjunction spaces.  Thus degree can control an entire family of CW homotopy types.
\end{solution}

% END VOL08-EXPANSION VIII08
